\documentclass[12pt]{article}

\usepackage[T1]{fontenc}
\usepackage[utf8]{inputenc}
\usepackage[brazil]{babel}
\usepackage{lmodern}
\usepackage{microtype}
\usepackage[a4paper,margin=2.5cm]{geometry}
\usepackage{setspace}
\usepackage{indentfirst}

\setstretch{1.15}

\title{As (Pelo Menos) Oito Leis de Say, ou o Que Say e James Mill Podem Realmente Ter Querido Dizer \thanks{Texto Traduzido livremente por Luiz Mario Andrade com o auxílio da ferramenta Chat GPT 5.4}}
\author{William J. Baumol}
\date{Maio de 1977}

\begin{document}

\maketitle

\begin{center}
\textit{Economica}, 44, 145-162\\
Universidades de Princeton e de Nova York
\end{center}

\section{A Discussão da Lei de Say como um Complexo de Ideias}

O revisionismo costuma fornecer toda a excitação que se pode encontrar nos escritos sobre história do pensamento econômico (\textit{dogmengeschichte}). Aqui, o êxito frequentemente consiste em mostrar que um autor foi mal interpretado, ou que sua contribuição é muito maior do que se supõe geralmente. Este artigo, entretanto, não pretende propor mais do que uma modesta revisão de nossas ideias sobre os primeiros escritos sobre a Lei de Say. Argumentar-se-á que J. B. Say e seu contemporâneo James Mill, a quem também se credita por vezes a noção de que "a oferta cria sua própria demanda", parecem ter tido em mente um conjunto de ideias consideravelmente mais complexo do que isso e, além disso, que as principais implicações de política econômica que extraíram de sua discussão foram muito além do pensamento reconfortante de que os temores de uma superabundância (\textit{glut}) universal são infundados.

Certamente, a discussão deles não foi de espírito panglossiano, pois um de seus pontos principais era a superioridade, do ponto de vista do crescimento do bem-estar econômico, dos gastos de investimento sobre o consumo ocioso (seja na forma de vida luxuosa por indivíduos ricos ou nos gastos militares e outros gastos improdutivos de governos pródigos). Além disso, esses autores sustentavam que apenas a produção pode criar poder de compra real, de modo que uma comunidade empobrecida não tem condições de fornecer demanda efetiva em abundância. A restrição na produção deve necessariamente limitar a demanda efetiva e certamente não é uma forma de promovê-la, particularmente no longo prazo.

Estas são ideias, notar-se-á, que a literatura moderna não rejeita, pelo menos não com a universalidade com que nega a noção de que a oferta cria sua própria demanda. De fato, são pensamentos que ainda podem ter alguma relevância, particularmente para a teoria e política de desenvolvimento. Além disso, para a época em que esses autores escreviam, essas ideias certamente não podem ser consideradas uma manifestação de conservadorismo. Pelo contrário, podem ser interpretadas como uma resposta aos defensores da aristocracia rural e de seus padrões de consumo, e como apoio ao espírito expansionista da empresa capitalista que acompanhou a jovem revolução industrial.

Obviamente, não há limite para a extensão em que se pode subdividir as noções que compõem a discussão de Say sobre mercados ou trocas (\textit{débouchés}), suas observações relacionadas sobre o consumo ou os materiais paralelos em Mill. Limitar-me-ei àquelas relações que desempenharam um papel substancial em discussões posteriores, particularmente na literatura recente, e àquelas observações que parecem claramente ter estado em primeiro plano na mente dos autores, conforme indicado pelo grau de ênfase dedicado a elas.

Mas antes de chegar a isso, convém recapitular o conteúdo da Lei de Say, tal como é entendida hoje em dia, e em contraste com a proposição principal defendida pelo próprio Say.

Quando se discute a Lei de Say, a literatura desde Ricardo parece geralmente referir-se a uma de duas proposições — uma afirmação forte, que Becker e eu denominamos em outro lugar como "Identidade de Say", e uma variante mais fraca, que rotulamos de "Igualdade de Say".

A primeira delas, a identidade, é a afirmação de que ninguém jamais deseja reter dinheiro por qualquer período significativo de tempo, de modo que, como resultado, cada oferta (oferta) de uma quantidade de bens constitui automaticamente uma demanda por um feixe de outros itens de igual valor de mercado. Assim, como em uma economia de escambo, a oferta deve criar automaticamente sua própria demanda e uma superprodução geral de bens e serviços é logicamente impossível.

A segunda forma, a igualdade, admite a possibilidade de períodos (breves) de desequilíbrio durante os quais a demanda total por bens pode ser inferior à oferta total, mas sustenta que existem forças equilibradoras confiáveis que devem em breve unir as duas.

Deve-se notar que ambas as proposições referem-se ao curto prazo — o período relevante para a análise de crises econômicas, recessões e desemprego. É uma das minhas principais teses que muitas das proposições centrais de Say são bastante distintas tanto da igualdade quanto da identidade e que, de fato, não é até a segunda edição do \textit{Traité} em 1814 que a estrutura completa de qualquer uma das ideias é detalhada no livro.

Acredito que as questões correlatas que Say enfatizou em sua primeira edição nos dão alguma ideia da ordem de prioridades em sua mente, e isso, bem como a quantidade de espaço e ênfase que lhes dedicou, sugere-me que algumas das proposições associadas podem ter tido um significado ainda maior para ele. Para ajudar o leitor a julgar essas questões por si mesmo, incluí neste artigo uma tradução das seções pertinentes desse livro, que acredito serem interessantes em si mesmas, e que são bastante breves e fáceis de seguir.

\section{Sobre a Prioridade na Formulação da Lei de Say}

Este artigo pode, incidentalmente, lançar alguma luz sobre um assunto mais secundário — a questão da prioridade na formulação da Lei de Say.

Na história das ideias, poucas coisas são mais temerárias do que a atribuição da paternidade de alguma proposição a um indivíduo específico. A atribuição é pouco mais que um ato esportivo que servirá, em última análise, como um desafio para outros, que certamente encontrarão, no devido tempo, um predecessor que antecipou a ideia com mais ou menos exatidão.

Como é bem sabido, a Lei de Say é uma noção cuja linhagem é mais antiga do que o nome sugere. Suas raízes nos escritos dos Fisiocratas foram investigadas com grande cuidado por Spengler (1945). Existem afirmações precursoras claras em \textit{A Riqueza das Nações} (ver Sowell, 1972, pp. 15-17, e Spengler, 1945, pp. 182-184). Também foi sugerido que James Mill antecipou Say na formulação do princípio. Embora a primeira edição do \textit{Traité} tenha aparecido em 1803, argumentou-se que em 1807 Mill, o pai, em seu \textit{Commerce Defended}, publicou algo muito mais próximo da "Lei de Say" como ela é hoje geralmente entendida. Spengler (1945) e Sowell (1972) mostram que esta conclusão provavelmente repousa sobre uma leitura superficial da primeira edição de Say. Os leitores podem ter sido induzidos ao erro porque grande parte da discussão pertinente aparece muito mais tarde no livro do que o capítulo sobre \textit{débouchés}, o capítulo no qual a maior parte deste material foi coletada em edições subsequentes. Winch (1966, p. 34) resolve a questão de forma ainda mais eficaz. Ele nos lembra, como Sowell fez, que Mill conhecia o \textit{Traité} quando escreveu \textit{Commerce Defended} e, de fato, o cita ali. Além disso, ele mostra que Mill creditou explicitamente a Say a ideia. Eu argumentarei, no entanto, que a questão é um pouco mais complexa. Embora seja verdade que a versão de Mill da discussão não difira tanto da de Say, sugerirei que nenhum deles ofereceu uma exposição explícita da Lei de Say até algum tempo depois. Apenas a segunda edição do \textit{Traité} de Say, em 1814, descreve explicitamente a lógica da Lei de Say como a interpretamos hoje em dia.

Passemos agora às questões correlatas que Say e Mill enfatizaram em suas discussões.

\section{A Produção como Fonte de Renda Disponível}

O título do capítulo focal, "Des Débouchés", é frequentemente traduzido como "sobre mercados". No entanto, deve-se ter em mente que o termo "mercado" é usado como em "fazer mercado", para denotar a disponibilidade de demanda efetiva, não como uma instituição (como um mercado físico) que facilita o processo de troca. Talvez uma tradução melhor de \textit{des débouchés} seja "sobre escoadouros para bens". Começo com as proposições que são o tema principal do extremamente breve capítulo sobre \textit{débouchés} na primeira edição do \textit{Traité}. Estas podem ser resumidas como:

\textit{Primeira Proposição de Say.} O poder de compra de uma comunidade (demanda efetiva) é limitado e igual à sua produção, porque a produção fornece os \textit{meios} pelos quais os produtos podem ser comprados. Além disso,

\textit{Segunda Proposição de Say.} Os gastos aumentam quando a produção aumenta.

Note que a primeira proposição trata do \textit{poder} de compra, não das compras reais. Ela nos diz, como Keynes fez, que a produção é a fonte da demanda efetiva — que a produção é poder de compra. Mas não diz que todo esse poder de compra será sempre usado para comprar bens. Em vez disso, essas afirmações afirmam, de fato, que a propensão marginal a consumir e investir é maior que zero e que a propensão média (geralmente) não é maior que a unidade. Isto é, dizem-nos que as pessoas que produzem mais estão em posição de consumir mais e geralmente o farão. Esses pontos, que a maioria de nós ainda aceitaria, são feitos de forma mais enfática no brevíssimo capítulo sobre \textit{débouchés} na primeira edição, que agora apresento em sua totalidade:

\begin{quote}
Capítulo 22 - Sobre Mercados

Todo produtor produz uma quantidade de um bem específico que excede consideravelmente seu próprio consumo. O agricultor colhe mais grãos do que o necessário para alimentar a si e sua família; o chapeleiro faz substancialmente mais chapéus do que produz para seu uso; o merceeiro atacadista manipula mais açúcar do que pode consumir. Cada um deles precisa de muitos outros produtos para viver confortavelmente. As trocas que realizam de seus próprios produtos pelos de outros constituem o que se chama de \textit{mercados} para esses produtos.

Nesta operação, o dinheiro serve aproximadamente ao mesmo papel que os cartazes e os panfletos em uma grande cidade, que facilitam o intercâmbio de pessoas que queiram fazer negócios umas com as outras. Durante o curso do ano, cada produtor lida com uma quantidade muito grande de dinheiro, mas, exceto por pequenos saldos de pouca importância, ao final do ano geralmente não resta em suas mãos mais dinheiro do que tinha no início. O que importa é o que ele compra com esse dinheiro, ou seja, os produtos de outros que trocou pelos seus, uma parte dos quais consumiu e uma parte economizou de acordo com suas necessidades, seus hábitos de poupança e o estado desta riqueza.

Confio que isso mostre que não é a abundância de dinheiro, mas a abundância de outros produtos em geral que facilita as vendas. Esta é uma das verdades mais importantes da economia política.

Imagine um indivíduo muito industrioso que tenha tudo o que precisa para produzir coisas: tanto habilidade quanto capital; se ele fosse a única pessoa industriosa em uma população que, à parte alguns alimentos grosseiros, não sabe como fazer nada; o que ele poderia fazer com seus produtos? Ele comprará a quantidade de alimento grosseiro necessária para satisfazer suas necessidades. Mas o que pode fazer com o resíduo? Nada. Mas se a produção do país começar a se multiplicar e a tornar-se mais variada, então todo o seu produto poderá encontrar um uso, ou seja, poderá ser trocado por coisas de que necessita ou por luxos adicionais que possa desfrutar, ou para a acumulação dos estoques que considere apropriados.

O que acabo de dizer sobre um único indivíduo industrioso pode ser dito igualmente de cem mil. Sua nação lhes oferecerá tanto mercado quanto ela puder pagar por objetos adicionais; e ela pode pagar por objetos adicionais em proporção às quantidades que produz. O dinheiro desempenha nada mais do que o papel de um conduto nesta dupla troca. Quando as trocas forem completadas, verificar-se-á que se pagou por produtos com produtos.

Consequentemente, quando uma nação tem uma quantidade grande demais de um tipo específico de produto, o meio de descartá-lo é criar bens de outra variedade. É quando não se pode mais produzir qualquer objeto trocável que a exportação se torna vantajosa. Esse é o caso quando ela serve como meio para comprar produtos que não podem ser fornecidos domesticamente, como frutas de outro clima. Mas as vendas mais lucrativas são aquelas que uma nação faz para si mesma; pois elas só podem ocorrer por causa dos dois valores produzidos: o que é vendido e o que é comprado. A exportação deve, portanto, ser considerada apenas como um suplemento e menos vantajosa do que o consumo doméstico. [Say, 1803, Vol. I, Livro 1, pp. 152-155]
\end{quote}

Exatamente o mesmo ponto é enfatizado por James Mill:

\begin{quote}
Nenhuma proposição, entretanto, em economia política parece ser mais certa do que esta que vou anunciar, por mais paradoxal que possa parecer à primeira vista; e, se for verdadeira, nenhuma, sem dúvida, pode ser considerada de maior importância. A produção de mercadorias cria, e é a única e universal causa que cria um mercado para as mercadorias produzidas. Consideremos apenas o que se entende por mercado. Entende-se por ele algo mais do que o fato de algo estar pronto para ser trocado pela mercadoria de que queremos nos desfazer? Quando as mercadorias são levadas ao mercado, o que se quer é alguém para comprar. Mas para comprar, é preciso ter com que pagar. É obviamente, portanto, os meios coletivos de pagamento existentes em toda a nação que constituem o mercado total da nação. Mas em que consistem os meios coletivos de pagamento de toda a nação? Não consistem eles na sua produção anual, na renda anual da massa geral de seus habitantes? Mas se o poder de compra de uma nação é medido exatamente por sua produção anual, como indubitavelmente é; quanto mais se aumenta a produção anual, mais por esse próprio ato se estende o mercado nacional, o poder de compra e as compras reais da nação. Qualquer que seja a quantidade adicional de bens, portanto, que seja criada em qualquer país a qualquer momento, um poder de compra adicional, exatamente equivalente, é criado no mesmo instante; de modo que uma nação nunca pode estar naturalmente superestocada nem com capital nem com mercadorias; pois a própria operação do capital cria um escoadouro (\textit{vent}) para seu produto. [Mill, 1807, pp. 81-82].
\end{quote}

Deve-se notar que tanto Mill quanto Say aparentemente interpretaram mal a lógica de suas próprias proposições. Da afirmação válida de que um alto nível de renda \textit{permite} um alto nível de demanda, eles parecem saltar para a conclusão de que a demanda será \textit{necessariamente} alta quando a produção for alta. Esta é uma questão à qual voltaremos mais tarde.

\section{Investimento vs. Consumo como Estimulantes da Riqueza}

Há um par de proposições correlatas que podem ter tido uma importância ainda maior, pelo menos para Say, a julgar pelo espaço que ele lhes dedicou e pela veemência com que as defendeu. Essas observações, que tratam da relação entre consumo e investimento (gastos "improdutivos" e "reprodutivos"), podem ser descritas como:

\textit{Terceira Proposição de Say.} Um determinado gasto de investimento é um estimulante muito mais eficaz para a riqueza de uma economia do que uma quantidade igual de consumo.

Esta proposição é a substância principal do capítulo de Say sobre o consumo, cuja tradução segue:

\begin{quote}
Capítulo 3 - A Riqueza de um Estado é Aumentada pelo seu Consumo?

O consumo reprodutivo que ordinariamente substitui por valores maiores aqueles valores que destrói, não é ordinariamente referido como \textit{consumo}, e eu mesmo, quando porventura usei essa palavra sem explicação, quis dizer consumo improdutivo, cujo único propósito é satisfazer as necessidades dos homens ou multiplicar seus prazeres. É apenas a esse tipo de consumo que se refere o título deste capítulo.

Muitas pessoas, vendo que, no geral, a produção total sempre iguala o consumo (porque é perfeitamente necessário que tudo o que é consumido tenha sido produzido), imaginaram que o incentivo ao consumo era favorável à produção. Os Fisiocratas agarraram-se a essa ideia e fizeram dela um dos princípios fundamentais de sua doutrina. O consumo é a medida da reprodução, diziam eles;\footnote{Ver Mercier de la Rivière, \textit{Ordre essentiel des Sociétés polit}, Vol. 2, p. 138, e os outros escritos dos Fisiocratas.} isto é, quanto mais se consome, mais se produz. E como a produção traz riqueza, concluíram que um Estado se enriquece através de seu consumo, que a poupança está em conflito direto com a prosperidade pública, e que o cidadão mais útil é aquele que mais gasta.\footnote{Um revisor observa que o conceito fisiocrata de "consumo" era muito diferente do de Say (ou do nosso), de modo que "a leitura de Say de sua própria definição nas palavras dos fisiocratas criou uma distorção completa de suas doutrinas." [W.B.]}

Este sistema é bem desenhado para ganhar o favor do vulgo e por isso tem muitos partidários. O fabricante e o comerciante percebem a prosperidade geral apenas na venda de suas mercadorias, apenas no maior consumo possível delas. Mas quando se considera esta doutrina mais atentamente, descobre-se que ela leva a resultados bastante diferentes.

O consumo de cada família pode exceder, igualar ou ser inferior à sua renda. O consumo de todas as famílias juntas, isto é, o da nação, pode seguir o mesmo curso. Ou seja, uma nação, compensadas todas as coisas, pode gastar mais ou menos do que sua renda ou exatamente o montante de sua renda.

No primeiro desses casos, a nação a cada ano causa um dano em seu capital; e, como resultado, a cada ano diminui sua renda; primeiro, pelos lucros que teriam correspondido ao capital que foi "comido" e, segundo, pelos lucros da indústria que esse capital teria sustentado. Longe de ser a forma de estimular as trocas, é a forma de reduzi-las. A cada ano a produção eleva-se ao nível do consumo; mas declina junto com o consumo, até o ponto em que a nação inteira, sem mais capital, sem mais terras cultivadas, sem mais indústria, sem mais população, desaparece inteiramente ou persegue uma existência triste e miserável. É a situação em que muitas partes da Grécia e da Síria caíram sob o domínio turco.

Se um país ou as famílias que o compõem não consomem mais do que sua renda, como esse país não esgota nenhuma parte de seu capital, ele manterá sua renda constante e oferecerá todos os anos o mesmo mercado para o produto de sua indústria.

Notemos também, de passagem, que é difícil para uma nação encorajada a gastar toda a sua renda não a gastar excessivamente rápido e, consequentemente, não cair mais ou menos rapidamente nos mesmos extremos penosos que aquelas que consomem uma parte de seu capital junto com sua renda. Tanto os indivíduos quanto um governo que elevam suas despesas regulares exatamente ao nível de sua renda ordinária não levam em conta acidentes e riscos imprevistos que nunca deixam de tirar algo da renda e de adicionar algo às despesas.

Quanto a uma nação que não gasta toda a sua renda, e anualmente aumenta seu capital, esta é a única que fornece os maiores mercados anuais para seu produto. Com efeito, a cada ano ela experimenta crescimento nos lucros de seu capital e no poder de sua indústria e, consequentemente, em sua renda; isto é, em seus meios de consumo, seja direto ou através de troca, em uma palavra, seus mercados.

O interesse público não é, consequentemente, servido pelo consumo, mas é servido prodigiosamente pela poupança, e embora pareça extraordinário para muitas pessoas, não sendo menos verdadeiro por consequência, a classe trabalhadora é servida por ela mais do que qualquer outra. Essas pessoas pensam, talvez, que os valores que os ricos economizam dos gastos em seus prazeres pessoais a fim de adicionar aos seus capitais não são consumidos. Eles \textit{são} consumidos; eles fornecem mercados para muitos produtores; mas são consumidos \textit{reprodutivamente} e fornecem mercados para os bens úteis que são capazes de engendrar ainda outros, em vez de serem evaporados em consumo frívolo.

Explicarei esta doutrina através de uma ilustração expressa em termos das atividades mais comuns.

Um indivíduo rico que tem uma renda de cem mil francos, e que tem o hábito de consumi-los totalmente, decide um dia diminuir seus gastos. Ele abre mão de alguns de seus servos domésticos, e é melhor servido; compra menos joias e não é tão criticado; oferece jantares menos esplêndidos e faz melhores amigos. Em resumo, em vez de gastar cem mil francos por ano, não gasta mais do que oitenta mil. A partir do primeiro ano, ele adiciona vinte mil francos ao seu capital produtivo. Os cem mil francos que constituem sua renda são sempre gastos em sua totalidade, mas não mais do que oitenta mil deles são gastos improdutivamente; os vinte mil restantes são gastos de uma maneira que os reproduz com lucro. Ele os empresta a um fabricante de lenços cuja empresa estava definhando por falta de fundos. A classe que compreende os lacaios, joalheiros e mercadores de alimentos finos experimenta de fato uma diminuição na demanda por seus serviços e seus produtos; mas a classe que fornece aos fabricantes de lenços suas roupas, seu alimento e suas matérias-primas vê a sua aumentar precisamente na mesma proporção. O incentivo dado pelo consumo reprodutivo é, portanto, o mesmo que teria resultado da satisfação das necessidades e prazeres de uma única pessoa. Mas não para por aqui.

Com efeito, houve \textit{adicionalmente} um aumento na renda para o capitalista rico e um aumento na renda para o fabricante e seus trabalhadores. Com base em uma avaliação muito moderada baseada na experiência,

\begin{center}
\begin{tabular}{ll}
o capitalista pode encontrar sua renda aumentada em & 1.000 francos \\
o fabricante principal pode encontrar a sua aumentada pelo mesmo montante & 1.000 francos \\
e a classe de trabalhadores pode encontrar seus salários totais aumentados em & 3.000 francos \\
\hline
total & 5.000 francos
\end{tabular}
\end{center}

O consumo daquele ano, consequentemente, será de 105 mil francos em vez dos 100 mil a que teriam somado se a pessoa com a grande renda tivesse gasto tudo improdutivamente.\footnote{Pode-se pensar que há uma dupla contagem aqui, e que o lucro do empresário e de seus empregados são meramente substitutos para aqueles que teriam sido ganhos pelo joalheiro e pelo chef na suposição de consumo improdutivo. Mas não há realmente dupla contagem. Os lucros do joalheiro, do chef e de todos os fornecedores de consumo improdutivo são substituídos pelos fornecedores de consumo produtivo. Estes fornecedores venderam ao fabricante ou aos seus trabalhadores 20 mil francos em bens, no lugar da mesma quantidade de bens que teriam sido vendidos ao consumidor rico. Além disso, houve 5 mil francos de renda pelos quais o consumo do ano pode ter aumentado. Isto é, com 20 mil francos e trabalho, 25 mil francos em lenços foram produzidos. Não se deve surpreender que eu avalie o produto bruto de um fabricante em 25 por cento de seu capital; o próprio produto líquido frequentemente atinge esta proporção.} E o que é ainda mais, os mesmos 20 mil francos que aumentaram as rendas daquele ano em 5 mil francos são reacumulados e podem render o mesmo serviço em todos os anos que se seguem, enquanto se considerar apropriado usá-los produtivamente.

Há mil maneiras de investir as economias. Enquanto os efeitos felizes que acabamos de ver continuam, o ano seguinte dá ao mesmo consumidor a oportunidade de realizar economias semelhantes. Ele as usa para a construção de uma máquina a vapor, para irrigar um campo; o efeito sobre a riqueza geral é o mesmo: o rendimento de suas terras aumenta em mil francos, mais ou menos, e a atividade que é introduzida em seu cultivo rende uma renda anual a muitos trabalhadores.

Assim, muitas pessoas estão em erro quando imaginam que os pobres obtêm recursos apenas dos gastos dos ricos. O verdadeiro recurso do homem pobre é o seu trabalho. Para exercer seu trabalho, ele não precisa do consumo do homem rico; ele precisa apenas de seu capital. Da mesma forma, um país ou um distrito poderia ser muito afortunado se nenhum homem rico residisse ou consumisse ali, desde que investissem seu capital nele. O agricultor trabalharia ali para o fabricante e o comerciante, o comerciante para o agricultor e o fabricante, e o último para os outros dois. Todos estariam bem supridos de todas as necessidades da vida. Com frugalidade eles poderiam enriquecer-se e teriam também os meios para pagar ao capitalista ausente os juros e os aluguéis sobre o capital e a terra que, por assim dizer, ele lhes teria emprestado.

Insisti em realizar esta demonstração porque o erro que ela expõe está entre os mais difundidos. É compartilhado por aqueles que defendem a doutrina mercantilista e por aqueles que defendem a doutrina agriculturista, isto é, os Fisiocratas. Todos eles consideram o consumo útil em relação à produção, enquanto ele é, de fato, útil apenas pelo prazer que proporciona.\footnote{Conheci um jovem que jogava frascos de cristal pela janela assim que os esvaziava. "Isso é necessário", dizia ele, "para incentivar as oficinas". Dessa forma, ele diminuía a riqueza social total, precisamente na proporção em que o artesão a aumentava ao produzir o frasco. Ele deveria tê-lo dado a uma família pobre demais para conseguir um. Dessa forma, ele teria proporcionado aos seus habitantes um dos prazeres da riqueza, e a manufatura teria recebido uma vantagem comparável àquela que teria recebido de um aumento na opulência geral.} Um consumidor não oferece vantagens por aquilo que consome, mas por aquilo que fornece como reposição; ora, ele pode dar muito mais como reposição se realizar menos consumo sem recompensa e mais consumo reprodutivo. Se um homem rico ocioso que devora uma renda enorme não arruína seu país, pelo menos não contribui em nada para sua prosperidade. Um homem rico que não gasta toda a sua renda contribui para ela; um homem rico que aumenta seu capital e, além disso, se ocupa utilmente, contribui ainda mais.

Se existe um hábito que merece encorajamento, tanto em monarquias quanto em repúblicas, em países grandes e pequenos, é a poupança. Mas ela precisa de encorajamento? Não basta evitar honrar a dissipação? Não basta respeitar todas as poupanças e todos os seus usos, isto é, o crescimento de toda atividade econômica que não seja criminosa? Pois não estou falando aqui da proteção vaga e inadequada que às vezes é referida como \textit{respeito à propriedade}. Não basta apenas proteger os indivíduos de serem espoliados por atos arbitrários, pela injustiça, pela fraude e violência; é também necessário garanti-los contra ataques por chicana, contra aumentos de impostos, mesmo que livremente consentidos, porque, embora uma contribuição possa ser votada livremente, ela se torna obrigatória no momento em que é transformada em requisito. Então, o desejo do homem de acumular bens, de preparar seus recursos para o futuro, bastará para equilibrar seu amor pelo prazer presente e para prover a prosperidade geral que é feita da prosperidade dos indivíduos. [Say, 1803, Vol. II, Livro 5, pp. 358-368]
\end{quote}

Embora nesta primeira edição do \textit{Traité} a questão da superioridade do investimento sobre o consumo como estímulo ao crescimento não seja referida diretamente no capítulo sobre \textit{débouchés}, ela é explicitamente incorporada a este capítulo em edições posteriores (ver, por exemplo, p. 139 da tradução inglesa da terceira edição). Para Mill, também, este é um ponto central: "A grandeza do produto de um país em qualquer ano depende inteiramente da grandeza da quantidade do produto do ano anterior que é destinada ao negócio da reprodução. O produto anual é, portanto, tanto maior quanto menor for a parte destinada ao consumo" (Mill, 1807, p. 71).

Mill cita com aprovação o conto de Say sobre o jovem que jogava vasos de cristal pela janela para incentivar as manufaturas. Ele também ressalta que

\begin{quote}
... é a manutenção de grandes frotas e exércitos que é sempre o peso mais formidável na balança do consumo, e que tem a tendência mais fatal de virar o equilíbrio contra a reprodução e a prosperidade. É pela lamentável artimanha das guerras, quase sempre alimentadas por preconceitos infantis e paixões cegas, que os negócios de nações prósperas são primeiro levados à condição estacionária, e a partir desta mergulhados no retrógrado. [Mill, 1807, p. 74]
\end{quote}

(Esta passagem também nos lembra das circunstâncias em que Mill estava escrevendo. Ele estava aqui envolvido em um debate sobre o comércio exterior, a hegemonia da agricultura, gastos de guerra e subconsumo — questões surgidas do bloqueio tentado por Napoleão às Ilhas Britânicas.)

\section{No Longo Prazo a Demanda Alcançará a Oferta}

Como todos os economistas clássicos, Say estava muito interessado no longo prazo. Desde a primeira edição, Say nunca deixou de enfatizar, como uma observação empírica:

\textit{Quarta Proposição de Say.} Ao longo dos séculos, a comunidade sempre encontrará demandas para produções aumentadas, mesmo para aumentos que sejam enormes.

Em outras palavras, não havia, em sua opinião, base para o temor da estagnação secular. Esta visão ele baseou no argumento de que a produção fornece o \textit{poder de compra} com o qual a produção pode ser adquirida. Assim, desde a primeira edição, ele recordava ao leitor quanto a produção de seu país havia aumentado desde os dias humilhantes de Henrique V e o fim da Guerra dos Cem Anos. Ele aponta que a renda total auferida pelos fatores de produção expande-se igualmente com a produção total da comunidade. "Caso contrário, por que meios se poderia comprar hoje em dia na França pelo menos duas ou três vezes a quantidade de coisas que eram compradas no miserável e infeliz reinado de Carlos VI?" (Say, 1803, Vol. II, p. 180.) Carlos VI (1368–1422) foi o rei mentalmente instável da França que foi forçado pelos ingleses a adotar Henrique V como seu herdeiro, deserdando seu filho, o futuro Carlos VII, e a dar a Henrique sua filha Catarina em casamento. Desenvolvimentos subsequentes parecem, em geral, ter apoiado Say neste ponto. O PIB per capita real sem dúvida cresceu pelo menos dez vezes desde que Say escreveu, quando a Revolução Industrial estava apenas no início. A demanda certamente parece ter crescido de forma comparável.

Ora, deve-se admitir que Say parece nunca ter dedicado muito espaço a uma discussão da correspondência de muito longo prazo entre demanda e oferta, e não encontrei nenhum lugar onde Mill tenha tratado disso. No entanto, parece-me estar implícito em sua ênfase na produção como a fonte de \textit{crescimento} na riqueza de uma nação. O fato de que, em sua segunda edição, Say moveu sua observação sobre Carlos VI para o corpo do capítulo sobre \textit{débouchés} e que a manteve ali nas edições subsequentes atesta o valor que ele dava à observação de que a demanda no longo prazo é capaz de acompanhar os aumentos mais enormes na produção.

Note que esta proposição, que é claramente de considerável importância para a análise do desenvolvimento econômico, é virtualmente irrelevante para a política de estabilização. A observação de Say de que no longo prazo a demanda acompanha os aumentos na produção é perfeitamente consistente com a existência de períodos prolongados de desemprego substancial. Não há nada no princípio que seja inconsistente com uma "superabundância geral" (\textit{general glut}) cuja possibilidade é negada pela identidade de Say. Note, a este respeito, a afirmação de Sowell (1974, pp. 64ss) (contra Patinkin e Blaug) de que mesmo a posição de Malthus na controvérsia da superabundância não incluía um argumento de estagnação secular.

\section{O Dinheiro Não é Mais do que um Meio para Facilitar a Troca}

Chegamos agora à primeira das proposições que se relacionam estreitamente com a "Lei de Say". Este é o destaque que tanto Say quanto Mill dão à falta de importância do dinheiro como contribuinte para a demanda efetiva. Say sublinha a conclusão de que o dinheiro nada mais é do que um conduto para, ou um lubrificante do, processo de troca, mas que em si mesmo, em última análise, não contribui em nada para a demanda efetiva. (Mill trata do papel do dinheiro apenas obliquamente. Ele diz que suas proposições são vistas mais claramente considerando "...as circunstâncias de um país no qual todas as trocas fossem por via de escambo, pois a ideia de dinheiro frequentemente tende a perplexar" (p. 82)). Ele também faz muito uso da observação de Spence: "Que não se alegue que o que eles pudessem economizar seria entesourado (pois os avarentos de hoje são mais sábios do que guardar seu dinheiro em caixas fortes em casa), mas seria emprestado a juros", enfatizando que esses empréstimos financiarão a demanda efetiva na forma de consumo reprodutivo (p. 75). Já vimos que dois dos sete parágrafos que constituem o capítulo da primeira edição sobre \textit{débouchés} são dedicados a uma descrição da irrelevância do dinheiro. Say estabelece este ponto comparando o papel do dinheiro ao da publicidade e enfatizando que os produtores geralmente adicionam quantias desprezíveis ao seu estoque de dinheiro ao longo de um ano. Assim, podemos descrever:

\textit{Sexta Proposição de Say.} A produção de bens, e não a oferta de dinheiro, é o determinante primário da demanda. O dinheiro facilita o comércio, mas não determina as quantidades de bens que são trocados.

Esta afirmação aproxima-se mais de implicar a identidade de Say do que qualquer outra coisa que discutimos até agora. Se de fato ele tivesse dito que ninguém quer dinheiro por si mesmo e que quaisquer acumulações de caixa serão sempre e imediatamente gastas, a identidade de Say teria de fato sido implicada. No entanto, para o Say da primeira edição, a irrelevância do dinheiro pareceria surgir de outra maneira — ele é irrelevante não porque apenas bens sejam demandados, mas porque é um meio pobre para o estímulo da demanda. Voltamos ao ponto keynesiano de que a demanda efetiva é uma função crescente da produção \textit{real}, e que o dinheiro apenas facilita o funcionamento dessa relação. Se é isso que Say tinha em mente (e sobre isso só podemos especular a partir de sua redação inconclusiva), então mesmo sua discussão monetária está apenas distantemente relacionada à "Lei de Say".

\section{Superabundâncias Gerais vs. Particulares}

Até onde pude julgar, Say em sua primeira edição e Mill em \textit{Commerce Defended} fazem apenas um ponto que está realmente relacionado diretamente à forma de igualdade ou identidade da Lei de Say. Esta é uma afirmação que podemos caracterizar como:

\textit{Sétima Proposição de Say.} Qualquer superabundância (\textit{glut}) no mercado para um bem deve envolver subprodução relativa de alguma outra mercadoria, ou mercadorias, e a mobilidade do capital para fora da área com excesso de oferta e para as indústrias cujos produtos são insuficientes para atender à demanda tenderá rapidamente a eliminar a superprodução.

Certamente este argumento sugere fortemente que ambos os autores acreditavam em alguma forma da Lei de Say. No entanto, este argumento parece ser uma conclusão que se tira da Lei, em vez de um fragmento de lógica que defenda sua validade. É um corolário, em vez de um fundamento da Lei.

Mill faz o ponto com força característica:

\begin{quote}
Pode ser necessário, entretanto, observar que uma nação pode facilmente ter mais do que o suficiente de qualquer mercadoria específica, embora nunca possa ter mais do que o suficiente de mercadorias em geral. A quantidade de qualquer mercadoria pode facilmente ser levada além de sua devida proporção; mas por essa própria circunstância está implícito que alguma outra mercadoria não é fornecida em proporção suficiente. O que de fato se entende por uma mercadoria exceder o mercado? Não é que haja uma parte dela para a qual nada há que possa ser obtido em troca? Mas daquelas outras coisas, então, a proporção é pequena demais. Uma parte dos meios de produção que foi aplicada na preparação desta mercadoria superabundante deveria ter sido aplicada na preparação daquelas outras mercadorias até que o equilíbrio entre elas tivesse sido estabelecido. Sempre que este equilíbrio é devidamente preservado, não pode haver superfluidade de mercadorias, nenhuma para a qual um mercado não esteja pronto. Este equilíbrio também a ordem natural das coisas tem uma tendência tão poderosa a produzir, que ele será sempre preservado com muita exatidão onde a interferência imprudente do governo não o impedir, ou aquelas desordens no intercâmbio do mundo, produzidas pelas guerras nas quais a parte inofensiva da humanidade é mergulhada, pela tolice com muito mais frequência do que pela sabedoria de seus governantes. [Mill, 1807, pp. 84-85]
\end{quote}

Note que nesta passagem Mill não apenas enfatiza a possibilidade de uma superabundância parcial, mas também afirma explicitamente a impossibilidade de um excesso "de mercadorias em geral". Em outras palavras, Mill afirma sua adesão à Lei de Say. No entanto, parece-me que em nenhum lugar aqui ou na primeira edição de Say qualquer um dos autores tenta detalhar o fundamento lógico para qualquer variante da Lei de Say. Para provar este ponto, apresentarei agora uma última passagem do \textit{Traité}. Esta é talvez particularmente significativa porque tanto Spengler (1945) quanto Sowell (1972) a citam como a passagem contendo o cerne da primeira discussão de Say sobre o assunto.

\begin{quote}
Capítulo 5 - Em Que Proporções o Valor dos Produtos é Distribuído Entre os Três Fatores de Produção

A magnitude da demanda por fatores de produção em geral não depende da magnitude do consumo como tantas pessoas imaginaram. O consumo não é de todo uma causa, é um efeito. Para consumir é necessário comprar; ora, pode-se fazer compras apenas com o que se produziu. A quantidade de produtos demandados é, consequentemente, determinada pela quantidade de produtos criados? Sem dúvida alguma. Todos podem, ao seu bel-prazer, consumir o que produziram; ou então podem comprar outro produto com o seu. A demanda por produtos em geral é, portanto, sempre igual à soma dos produtos disponíveis.\footnote{Um produto consumido por seu fabricante é um produto que foi ofertado e demandado pela mesma pessoa. Ele constitui parte da oferta total de bens e da demanda total por eles.} Uma nação que produzisse anualmente não mais do que um valor de dois bilhões não seria capaz de comprar nem de consumir três bilhões durante o mesmo período sem retirar o déficit de um bilhão de seus capitais a cada ano. Vemos que a melhor maneira de prover mercados para a produção é produzir mais, não destruí-los. Se esta conclusão é óbvia, como acredito ser, o que devemos pensar dos sistemas que incentivam o consumo para estimular a produção?

Não é verdade que a produção total não pode exceder o consumo total. Não se pode acumular uma parte dos produtos criados a cada ano? Não pode cada um acumular alguns de seus próprios produtos ou alguns daqueles que adquire por troca? Não é um mercado provido por esta acumulação tão eficazmente como se esse mesmo valor tivesse sido consumido?\footnote{Além disso, veremos mais tarde que a parte economizada da renda e adicionada aos capitais é tão certamente consumida a cada ano, mas de outra maneira; isto é, de uma maneira reprodutiva.} A produção total não é, consequentemente, limitada pela magnitude do consumo. A limitação do consumo não fecha os mercados; pelo contrário, o fomento da produção abrirá novos mercados. Uma nação que se enriquece desfruta de uma vantagem comparável àquela adquirida por uma que estende seu comércio exterior. Ela experimenta a abertura de novos mercados e o aparecimento de novos clientes. Ela estende seu comércio e não trava guerras para conseguir isso.

Se a produção não é limitada pela magnitude do consumo,\footnote{Aqui nosso assunto é apenas o consumo improdutivo, não o uso do capital para sua reprodução com lucro.} se se pode produzir mais do que se consome, quais são então os limites da produção? Eles consistem na disponibilidade dos fatores de produção.

Mas, dir-se-á, se há bens que não podem ser vendidos, há necessariamente mais fatores produtivos empregados do que oportunidades para o consumo de sua produção. De forma alguma. Nenhuma superabundância jamais ocorre, exceto quando uma quantidade grande demais de fatores de produção é dedicada a um tipo de produção e não o suficiente a outro. Com efeito, qual é a causa da incapacidade de realizar uma venda? É a dificuldade de obter algum outro bem (seja um produto ou dinheiro) em troca do que é oferecido. Meios de produção estão, consequentemente, faltando para o primeiro na medida em que são superabundantes para o segundo. Uma região no interior profundo de uma terra não encontra venda para seu trigo, mas se alguma fábrica for estabelecida ali e parte do capital e do trabalho que anteriormente era dedicado à terra for redirecionada para outro tipo de produção, os produtos de um e de outro podem ser trocados sem dificuldade, embora essas produções tenham se expandido em vez de diminuir. A incapacidade de vender, portanto, surge não da superabundância, mas da má alocação dos fatores de produção.

Nas notas que Garnier incluiu em sua excelente tradução de Smith, ele afirma que em nações antigas como as da Europa, nas quais o capital foi acumulado ao longo de vários séculos, a superabundância do produto anual obstruiria o comércio, se não fosse absorvida por quantidades proporcionais de consumo. Percebo que o comércio pode ser obstruído pela superabundância de produtos específicos. É um mal que nunca pode ser nada além de temporário, pois a participação na produção de bens cujas saídas excedem a necessidade por eles e cujo valor é depreciado cessará rapidamente e, em vez disso, será dedicada à produção dos bens que são procurados. Mas não posso conceber que os produtos do trabalho de uma nação inteira possam jamais ser superabundantes, uma vez que um bem fornece os meios para comprar o outro. A soma de suas produções compõe a riqueza total de uma nação, e a riqueza não é mais um embaraço para as nações do que é para os indivíduos.

Tendo este ponto sido minuciosamente esclarecido, ele nos fornece uma resposta para a questão que nos ocupa e que repito: de que elementos depende a demanda por fatores de produção em geral? Ela depende do volume da produção e, como o volume da produção depende da quantidade de fatores de produção, a demanda por fatores de produção expande-se proporcionalmente à quantidade de meios de produção em si mesmos. Isto é, como resultado, uma nação sempre tem os meios para comprar tudo o que produz. Caso contrário, por que meios se poderia comprar hoje em dia na França pelo menos duas ou três vezes a quantidade de coisas que eram compradas no miserável e infeliz reinado de Carlos VI? [Say, 1803, Vol. II, Livro 4, pp. 175-180]
\end{quote}

Esta passagem resume claramente todas as sete proposições que listamos até agora. Em particular, afirma enfaticamente que "...o comércio pode ser obstruído pela superabundância de produtos específicos... bens cujas saídas excedem a necessidade por eles... Mas... os produtos do trabalho de uma nação inteira nunca podem ser superabundantes porque um bem fornece os meios para comprar o outro".

\section{A Incompletude da Enunciação da Lei}

O que falta, então, nesta passagem para qualificá-la como uma enunciação explícita da Lei de Say? A Lei é certamente enunciada na última parte da afirmação de Mill "Que uma nação pode facilmente ter mais do que o suficiente de qualquer mercadoria específica, \textit{embora nunca possa ter mais do que o suficiente de mercadorias em geral}" (Mill, 1807, p. 84, meus itálicos); e o mesmo é obviamente verdade para Say quando nos diz que "Os produtos do trabalho de uma nação inteira nunca podem ser superabundantes".

Assim, a proclamação da Lei de Say é encontrada em ambos os escritos. Como já sugeri, tudo o que falta é a sua lógica. Assim, note que Say completa a citação anterior dizendo-nos que "...não pode conceber que os produtos do trabalho de uma nação inteira possam jamais ser superabundantes visto que um bem fornece os \textit{meios} para comprar o outro" (meus itálicos). Em outras palavras, a comunidade \textit{pode pagar} para comprar os bens \textit{se desejar} fazê-lo. Voltamos ao que rotulamos como a primeira proposição de Say, que afirma que a produção cria o \textit{poder} de comprar a si mesma.

Mas isso evita a questão básica em toda a controvérsia subsequente sobre a Lei de Say. A questão não é se o poder de compra necessário para comprar a produção de uma nação está disponível, mas se ele será, de fato, usado.

Mill trata o assunto virtualmente como uma tautologia e, portanto, como algo cuja lógica não exige discussão: "O que de fato se entende por uma mercadoria exceder o mercado? Não é que haja uma parte dela para a qual nada há que possa ser obtido em troca? Mas daquelas outras coisas, então, a proporção é pequena demais". (Mill, 1807, p. 85). Em outras palavras, todo excesso de oferta é, por definição, um excesso de demanda por outra coisa. Mas o que é esse algo mais? Deve ser necessariamente uma mercadoria?

\section{A Segunda Edição de Say: O Ingrediente Faltante Fornecido}

Existem, parece-me, dois ingredientes faltantes na discussão da Lei de Say nos parágrafos precedentes. Pois eles falham em discutir a natureza da demanda gerada pelos detentores do poder de compra da comunidade, e não lidam com o \textit{tempo} dessa demanda.

Outras discussões, tanto depois como antes, trataram de ambos os assuntos de forma muito explícita. Elas afirmaram que as pessoas sempre usarão seu poder de compra para buscar bens de consumo ou de capital em troca; isto é, demandarão bens em vez de saldos de caixa, porque agir de outra forma seria irracional. Além disso, afirmava-se (às vezes, mas nem sempre) que esta demanda deve acompanhar a aquisição de poder de compra através da produção \textit{virtualmente sem atraso}. Estas afirmações somam-se claramente à identidade de Say, com a oferta criando automática e imediatamente uma demanda equivalente, ou a uma forma forte da igualdade, com forças de equilíbrio rápidas e poderosas que, geralmente ou sempre, canalizam as demandas para mercadorias.

É curioso que Adam Smith já tivesse discutido ambos os assuntos quase trinta anos antes de Say. Preocupado em apoiar a ideia de que a poupança não reduz a demanda, ele trata das questões que acabamos de levantar em duas passagens bem conhecidas. A primeira examina a motivação da demanda:

\begin{quote}
Em todos os países onde há uma segurança tolerável, todo homem de senso comum procurará empregar todo o estoque que puder comandar, na busca de prazer presente ou lucro futuro... Um homem deve ser perfeitamente louco se, onde há segurança tolerável, não emprega todo o estoque que comanda, seja ele seu próprio ou emprestado de outras pessoas, em uma ou outra dessas... formas. [Smith, 1776, p. 268]
\end{quote}

Na segunda passagem, ele considera o tempo desses gastos:

\begin{quote}
Aquela parte de sua renda que um homem rico gasta anualmente é, na maioria dos casos, consumida por hóspedes ociosos e servos servis, que nada deixam para trás em troca de seu consumo. Aquela parte que ele economiza anualmente, como visando ao lucro, é \textit{imediatamente} empregada como capital, é consumida da mesma maneira, e \textit{quase no mesmo tempo também}, mas por um conjunto diferente de pessoas. [Smith, 1776, p. 321, meus itálicos]
\end{quote}

Em outras palavras, como já foi apontado antes, Smith quase enunciou a Lei de Say e certamente detalhou sua lógica de forma mais completa do que Say ou Mill em seus primeiros escritos sobre o assunto. As únicas lacunas na discussão de Smith são sua estipulação de que um estado de "segurança tolerável" deve prevalecer para que o investimento (\textit{ex ante}) seja sempre igual à poupança e sua falha em declarar explicitamente a conclusão de que a oferta total deve, consequentemente, ser sempre igual à demanda.

Say só chega a esses assuntos em sua segunda edição, onde nos diz em seu capítulo revisado e expandido sobre \textit{débouchés} que:

\begin{quote}
... Vale a pena notar que um produto, tão logo criado, oferece \textit{a partir daquele instante} um mercado para outros produtos em toda a extensão de seu próprio valor. Pois todo produto é criado apenas para ser consumido, seja produtivamente ou improdutivamente, e de fato para ser consumido \textit{o mais rapidamente possível}, uma vez que cada valor cuja realização é atrasada causa uma perda ao indivíduo que é atualmente seu possuidor do ganho de juros correspondente a esse atraso... Um produto é, portanto, tanto quanto todos possam organizar, destinado ao consumo mais rápido. A partir do momento em que existe, ele busca consequentemente outro produto com o qual possa ser trocado. O ouro e a prata não são exceção, pois logo que o comerciante faz uma venda, ele procura empregar o produto de sua venda. [Say, 1814, pp. 147-148, itálicos dele]
\end{quote}

A partir da segunda edição, Say também trata de forma muito mais direta o papel do dinheiro e sua irrelevância como determinante da demanda. Por exemplo, na quarta edição do \textit{Traité} lemos:

\begin{quote}
Se um comerciante disser: "Não quero outros produtos pelos meus lanifícios, quero dinheiro", ... Você diz que quer apenas dinheiro; eu digo que você quer outras mercadorias, e não dinheiro. Pois para que, de fato, você quer o dinheiro? Não é para a compra de matérias-primas ou estoque para seu comércio, ou víveres para seu sustento? ... Assim, dizer que as vendas estão paradas devido à escassez de dinheiro é confundir os meios com a causa...

Há sempre dinheiro suficiente para conduzir a circulação e o intercâmbio mútuo de outros valores, quando esses valores realmente existem. Caso o aumento do tráfego exija mais dinheiro para facilitá-lo, a carência é facilmente suprida... Em tais casos, os comerciantes sabem bem como encontrar substitutos para o produto que serve como meio de troca ou dinheiro;\footnote{Por letras à vista, ou pós-datadas, notas bancárias, créditos rotativos, baixas contábeis, etc. em Londres e Amsterdã.} e o próprio dinheiro logo aflui, por esta razão, que toda produção gravita naturalmente para aquele lugar onde é mais demandada... [Say, 1821, pp. 133-134]
\end{quote}

O que é notável neste comentário é a adesão implícita de Say, nesta passagem, a uma forma de igualdade, e não de identidade da lei. Ele admite que um excesso de demanda por dinheiro pode surgir, mas nos diz que será apenas temporário. Se for apenas local, suprimentos de dinheiro logo fluirão de outros lugares. E, em qualquer caso, o excesso de oferta será eliminado não por algo como o efeito saldo de caixa, que reduz as demandas para a oferta disponível (nominal) de caixa, mas por um aumento na oferta de caixa ou de substitutos de caixa.

Assim, a oitava (e para nossos propósitos a última) das oito proposições de Say é a própria Lei de Say. Aparentemente, isso assume a forma de um tipo de igualdade de Say, ou seja, a oferta e a demanda são sempre equalizadas por um mecanismo de equilíbrio rápido e poderoso. No entanto, devemos ter o cuidado de não atribuir a ele pensamentos sobre o que ele teria considerado distinções sutis, como aquela entre as formas de igualdade e identidade da lei. Pois, muito provavelmente, o assunto nunca foi explicitamente considerado por ele.

\section{Comentário Final}

Várias conclusões emergem desta discussão sobre a abordagem de Say e Mill à Lei de Say. Primeiro, na questão relativamente sem importância da prioridade, as coisas são nebulosas, e a resposta deve depender de nossa definição da própria Lei. Obviamente, ela é ao menos insinuada desde a primeira edição do \textit{Traité} em diante. Minha própria conclusão é que Spengler (1945), Winch (1966) e Sowell (1972) estão certos ao atribuir prioridade a Say. Mas eles estão certos pela razão errada. Para mim, a primeira enunciação da Lei, tal como a entendemos hoje, aparece completa com seus fundamentos lógicos nem na primeira edição do \textit{Traité} nem em \textit{Commerce Defended}, mas na segunda edição da obra de Say. Nessa interpretação, a Lei de Say pode ser dita como tendo alcançado sua primeira codificação completa em 1814.

No entanto, a principal conclusão que emerge do nosso reexame dos textos de Say e Mill é que a discussão da Lei de Say foi, antes de tudo, um exame das influências que promovem o crescimento econômico a longo prazo, e não primariamente uma questão de problemas de curto prazo de desemprego e superprodução. A ênfase principal dos argumentos de Say e Mill era que o investimento (consumo produtivo), em vez do consumo de luxo, construção de pirâmides ou gastos militares (consumo improdutivo), são os meios eficazes para promover o crescimento.

Há também uma segunda observação, indo além do nosso exame textual, que também parece valer a pena fazer. Tendemos a pensar hoje no lado malthusiano do argumento como sendo keynesiano em espírito e, consequentemente, "progressista", enquanto a posição Say-Mill-Ricardo é interpretada como sendo o oposto. Mas, vista em termos das circunstâncias do início do século XIX, quando o debate estava em curso, isso é o inverso da verdade. Malthus era o defensor aberto do papel dos proprietários de terras contra a classe capitalista em ascensão. (Marx caracteriza os escritos que "...consideram o consumo como um estímulo necessário para a produção" como "apologéticas... em parte para os ociosos ricos e os 'trabalhadores improdutivos' cujos serviços consomem, em parte para 'governos fortes' cujas despesas são pesadas, para o aumento das dívidas estatais, para detentores de sinecuras, etc." (1963, p. 281)). Foi em nome deles que ele defendeu tarifas protetoras e foi o consumo de bens de luxo e o emprego de retentores pessoais que ele estava defendendo quando argumentou que eles eram necessários para a prevenção da superprodução. Por outro lado, Say, Mill e Ricardo buscaram, contra esta posição, incentivar a nova atividade industrial e, com ela, a expansão da produção da economia. É sempre perigoso interpretar a cor política de obras anteriores em termos da situação da própria era do leitor.

\section*{Agradecimentos}

Sou grato a Elizabeth Bailey, Fritz Machlup e Donald Winch por seus comentários e sugestões, e a George de Menil e Arthur Laurie pelas correções de minhas traduções.

\section*{Referências}

\noindent BLAUG, MARK (1958). \textit{Ricardian Economics: A Historical Study}. New Haven: Yale University Press.

\noindent MARX, KARL (edição de 1963). \textit{Theories of Surplus Value}. Moscou: Progress Publishers, I.

\noindent MILL, JAMES (1807). \textit{Commerce Defended}. Londres: C. and B. Baldwin.

\noindent SAY, J. B. (primeira ed. 1803). \textit{Traité d'économie politique ou simple exposition de la manière dont se forment, se distribuent et se consomment les richesses}. Paris: Deterville.

\noindent ------ (2ª ed., 1814). \textit{Traité d'économie politique ou simple exposition de la manière dont se forment, se distribuent et se consomment les richesses}. Paris: A. A. Renouard.

\noindent ------ (1821). \textit{A Treatise on Political Economy; or the Production, Distribution and Consumption of Wealth}. Boston: Wells and Lilly.

\noindent ------ (1844). \textit{Cours complet d'économie politique pratique}. Bruxelas: Société Typographique Belge.

\noindent SMITH, ADAM (1776). \textit{An Inquiry into the Nature and Causes of the Wealth of Nations}. Londres: Methuen \& Co. (ed. Cannan 1925).

\noindent SOWELL, THOMAS (1972). \textit{Say's Law: An Historical Analysis}. Princeton: University Press.

\noindent ------ (1974). \textit{Classical Economics Reconsidered}. Princeton: University Press.

\noindent SPENGLER, J. J. (1945). The Physiocrats and Say's Law of Markets. \textit{Journal of Political Economy}, 53, 193-211, 317-347. Republicado em Spengler e Allen (1960).

\noindent ------ e ALLEN, W. R. (1960). \textit{Essays in Economic Thought: Aristotle to Marshall}. Chicago: Rand McNally.

\noindent TUCKER, G. S. L. (1960). \textit{Progress and Profits in British Economic Thought, 1650-1850}. Cambridge: University Press.

\noindent WINCH, DONALD (1966). \textit{James Mill: Selected Economic Writings}. Chicago: University Press.

\end{document}